% chapters/09/latihan.tex
\bagian{Latihan dan Refleksi}

\begin{enumerate}
  \item Jelaskan mengapa variabel bernilai benar-salah pada Bab~VIII sudah
        merupakan mesin keadaan yang paling sederhana, dan apa yang belum ada
        padanya.

  \item Sebutkan tiga lambang yang dipakai pada diagram keadaan--transisi
        beserta makna masing-masing.

  \item Gambarlah diagram keadaan--transisi untuk NPC yang memiliki keadaan
        berdiam, berpatroli, dan mengejar, dengan syarat yang dapat diperiksa
        pada setiap perpindahannya.

  \item Jelaskan tiga kesalahan yang sering terjadi pada diagram keadaan, dan
        gejala yang muncul pada program untuk masing-masing kesalahan.

  \item Sebutkan tiga jenis pewaktu yang dipakai pada bab ini beserta keperluan
        masing-masing.

  \item Jelaskan mengapa fungsi \code{decide} sengaja tidak mengubah keadaan
        apa pun. Apa manfaatnya bagi pengujian?

  \item Sebuah NPC berada pada keadaan menyerang selama 4,5 detik. Apabila
        pewaktu jeda serangannya 0,6 detik dan serangan pertama terjadi segera
        setelah keadaan dimasuki, hitunglah banyaknya serangan yang terjadi.

  \item Pada skenario pengujian, pemain menghilang pada jarak 300 piksel dari
        NPC, sedangkan jari-jari persepsi 260 piksel, kecepatan NPC 150 piksel
        per detik, dan pewaktu tenggang 1 detik. Tentukan apakah NPC benar-benar
        kehilangan pemain, beserta perhitungannya.
\end{enumerate}

\begin{refleksibox}
\begin{enumerate}
  \item Pada langkah kelima praktikum, keadaan menghindar yang Anda tambahkan
        memiliki syarat masuk dan keluar. Apakah kedua syarat itu dapat benar
        bersamaan? Bagaimana Anda memastikan tidak demikian?
  \item Ketika pemeriksaan keabsahan melaporkan perpindahan yang tidak sah pada
        langkah keenam, apakah laporannya mudah dipahami? Menurut Anda, apa
        yang membuat pesan pemeriksaan mudah atau sulit dipahami?
  \item Bandingkan NPC pada bab ini dengan NPC pada Bab~VIII. Menurut Anda,
        apakah perbedaan perilakunya sepadan dengan pertambahan panjang kodenya?
\end{enumerate}
\end{refleksibox}
